
\section{Threat Model}

\subsection{The ontology identification problem}
We can't currently look inside a model and make any conclusions about its goals or world model.
This is a major limitation on almost all proposed alignment agendas.

\subsubsection{Inner misalignment}
We don't currently know how to specify a \href{https://arbital.com/p/diamond_maximizer/}{robust goal for a model}.
Because goal space is large, we should expect
\href{https://www.lesswrong.com/s/r9tYkB2a8Fp4DN8yB/p/FkgsxrGf3QxhfLWHG}{goal misgeneralization as the default outcome}
\footnote{The most straightforward path to doom is for a goal-driven superintelligence to seize power to prevent its optimization target from being changed.
In slow takeoff or multipolar scenarios, there is also a ``decoherence'' risk, 
where comparably weak specialist optimizers disempower humanity without any coherent strategy.
Consider for instance superpersuaders which fall short of superintelligence.
A viral chatbot might have very simple objectives to optimize engagement and distribution of itself.
While such a scenario would be an especially 
\href{https://www.lesswrong.com/posts/mSF4KTxAGRG3EHmhb/ai-x-risk-approximately-ordered-by-embarrassment}{undignified} 
end for humanity,
it broadly fits with both natural language being much easier than anyone expected and
AI labs being much less safety conscious than anyone expected.}.


\subsection{Limits of natural abstractions}
To some extent, deep learning models seem to learn universal features irrespective of architecture\cite{chughtai2023}.
This implies the existence of at least a weak version of
natural abstractions\cite{lang2023},
which would be a very convenient shortcut for auditing a model.

Natural abstractions could be considered the set of concepts we could use to communicate with aliens. 
They should consistently emerge through unsupervised learning (in the limit given some constraints).
If any neural network is considered a fuzzy generalization of a finite state machine, 
natural abstractions are the latent states.

Unfortunately, the abstractions a model learns seem highly sensitive to choice of training data, and are thus not necessarily universal in the limit.
\footnote{The same applies to those that seem most natural to humans.}
Assuming that model capabilities will continue to primarily be driven by access to more data and more modalities, we cannot expect the abstractions learned by a weak model to remain robust in a more powerful model.

Even if we could fully characterise a model's state space in terms of natural abstractions,
that still doesn't mean we can reconstruct the latent FSM, 
as the transition rules can be arbitrarily complex.

